\documentclass[12pt]{article}
\usepackage{amssymb,amsmath,enumerate,graphics,graphicx} 


\newcommand{\cA}{{\cal{A}}}
\newcommand{\cU}{{\cal{U}}}
\newcommand{\D}{{\mathbb D}}
\newcommand{\ds}{\displaystyle}

\newtheorem{theorem}{Theorem}[section]
\newtheorem{lemma}[theorem]{Lemma}

{ \title{\sc A Verification of Campbell's Conjecture on Majorization-Subordination Theorems}}
\author{Roger W. Barnard \and Kent Pearce}
\date{30 July 2008}
\begin{document}

\maketitle
\begin{abstract}
In the early 70's, D.M. Campbell published three papers on majorization-subordination results for locally univalent functions.  
He showed: 

Theorem 1. \emph{ If $F$ is locally univalent of order $\alpha$ for $\alpha \ge 1$ and if $f$ is majorized by $F$ on 
$\{z : |z| < 1\}$, then $f'$ is majorized by $F'$ on $\{ z : |z| < n(\alpha) \}$ where 
$n(\alpha) = \frac{(\alpha+1)^{1/\alpha}-1}{(\alpha+1)^{1/\alpha}+1}$}.  

He further showed:  

Theorem 2. \emph{ If $F$ is locally univalent of order $\alpha$ for $\alpha \ge 1.65$ and if $f$ is subordinate to $F$ on 
$\{z : |z| < 1\}$, then $f'$ is majorized by $F'$ on $\{ z : |z| < m(\alpha) \}$ where 
$m(\alpha) = \alpha +1 - \sqrt{\alpha^2+2\alpha}$.}  

The bounds given in Theorem 1 and Theorem 2 are sharp for each $\alpha$.  
Campbell's results generalized similar majorization-subordination results which had been previously established 
assuming that the majorant was univalent.

The restriction $\alpha \ge 1.65$ in Theorem 2 arose out of the proof that Campbell constructed.  He conjectured that, 
in fact, Theorem 2 also held for $1.65 > \alpha \ge 1$.  We will review Campbell's proof of Theorem 2 and why the restriction 
$\alpha \ge 1.65$ arose in the proof.  We will then affirmatively verify Campbell's conjecture.

\vspace{0.4in}

\noindent Key words: majorization, subordination, locally univalent, Campbell \\
\noindent AMS Subject Classification 2000: 30C70

\end{abstract}

\section{Introduction}

Let $\cA(\D)$ denote the set of functions which are analytic on the open
unit disk $\D = \{z : |z| < 1\}$.  A function $\phi \in \cA(\D)$
is a \emph{Schwarz function} if $|\phi(z)| \le |z|$ for $z \in \D$.  
Let $f, F \in \cA(\D)$.  For $0 < r \le 1$, the function $f$ is \emph{majorized} 
by $F$ on $|z| < r$ if $|f(z)| \le |F(z)|$ on $|z| < r$.  The function 
$f$ is \emph{subordinate} to $F$ on $\D$, if there 
exists a Schwarz function $\phi$ such that  $f = F \circ \phi$ -- we write $f \prec F$.

Let $\cU_\alpha$ denote the subset of $\cA(\D)$ of all locally univalent
functions of order $\alpha$ with the normalization $f(z) = z+ \cdots$.
The family $\cU_\alpha$ is known as the universal linear invariant family of 
order $\alpha$ \cite{Pom01}.  In the early 70's Campbell published three papers 
on majorization-subordination for locally univalent functions \cite{DMC01,DMC02,DMC03}
in which he established relationships between majorization and subordination for functions
in $\cU_\alpha$.  In particular, he showed that many of the classical results on 
majorization and subordination for univalent functions hold in the more general setting 
of locally univalent functions.  He remarked ``Our investigation shows that 
the important datum for majorization-subordination theory, is not 
\emph{univalence}, but the order of a linearly invariant family."


Let $f,F \in \cA(\D)$ such that $F \in \cU_\alpha$ for some $\alpha \ge 1$.
Suppose $f \prec F$.  Then, there exists a Schwarz function $\phi \in \cA(\D)$ such that 
$f=F\circ\phi$. Let $a = f'(0)$.  We can 
assume that $f$ has been rotated so that $0 \le a \le 1$.  Then, we can write
${\ds \phi(z)=z\frac{a+\omega(z)}{1+a\omega(z)}}$ where $\omega \in \cA(\D)$ such that 
$\omega$ is a Schwarz function.

We will employ the following notation throughout the paper: For $\alpha \ge 1$, let 
$m(\alpha) = \alpha+1 - \sqrt{\alpha^2 + 2\alpha}$. For $z = xe^{it} \in \D$ 
with $0 \le x \le m(\alpha)$, let $c= \omega(z) = re^{i\theta}$.  We have 
$0 \le r \le x \le m(\alpha)$.  Note, at $z$ we have $f(z)=F(\phi(z))$ and 
${\ds \phi(z)= z\frac{a+c}{1+ac}}$.  


In \cite{DMC03}, Campbell established two lemmas (Lemma 1 and Lemma 2) for estimating $|\phi'(z)|$ in terms of $a$ and $c$, 
depending on whether $a$ was ``small" or ``large".  Then, referencing Pommerenke's classical paper \cite{Pom01}
on linearly-invariant families, he established the following fundamental inequality

\begin{equation} \label{fund}
\left | \frac{f'(z)}{F'(z)}\right | \le \frac{1-x^2}{1-|\phi(z)|^2}
\left( \frac{|1-\overline{\phi(z)}z|+|\phi(z)-z|}{|1-\overline{\phi(z)}z|-|\phi(z)-z|} \right ) ^\alpha
|\phi'(z)|
\end{equation}

Using his Lemma 2 \cite{DMC03}, Campbell showed that

\begin{equation} \label{fund-2}
\left | \frac{f'(z)}{F'(z)} \right | \le \frac{ba+1}{b+a} \left ( \frac{b-a}{b-1} \right ) ^ \alpha = k(a,\alpha,b)
\end{equation}

\noindent where ${\ds b = \frac{1+x^2}{2x} \ge \alpha+1}$, since $x \le m(\alpha)$. 
In particular, Campbell showed that $k(a,\alpha,b)$ is decreasing in $b$.  Hence, $k(a,\alpha,b) \le k(a,\alpha,\alpha+1)$. 
Following Campbell, we note that 

$$\frac{\partial k(a,\alpha,\alpha+1)}{\partial a} = \frac{(\alpha+1-a)^{\alpha-1}}{(\alpha)^{\alpha}}\frac{p(a,\alpha)}{(\alpha+1+a)^2}
$$

\noindent where 

\begin{equation} \label{poly_p}
p(a,\alpha) = -\alpha(\alpha+1)a^2 -(\alpha^3+3\alpha^2+4\alpha)a +\alpha(\alpha+1)^2. 
\end{equation}

\noindent Let $R=\{(a,\alpha) : p(a,\alpha) \ge 0 \}$, i.e., $R$ is the set where $k(a,\alpha,\alpha+1)$ is increasing
in $a$.  Campbell showed that $R$ contains the set $\{(a,\alpha) : 0 \le a \le 0.4, \alpha \ge 1 \}$.  

Let $C_1 = \{(a,\alpha) \in R : k(a,\alpha ,\alpha +1) = 1$ and $A_1 = \{(a,\alpha) \in R : k(a,\alpha ,\alpha +1) \le 1$. 
For $(a,\alpha) \in A_1$, we have $|f'(z)/F'(z)| \le 1$.  In this notation, Campbell explicitly verified that the rectangles
 
\begin{eqnarray*}
R_1 = \{(a,\alpha) : 0 \le a \le 3/20, 1.65 \le \alpha \le 2 \}\\
R_2 = \{(a,\alpha) : 0 \le a \le 1/6, 2 \le \alpha \le 3 \}\\ 
R_3 = \{(a,\alpha) : 0 \le a \le 1/10, 3 \le \alpha < \infty \}
\end{eqnarray*} 

\noindent each belong to $A_1$. More explicitly, Campbell showed that the ``right-hand" edges of $R_1, R_2$ and $R_3$ belong to $A_1$.  Figure 1 summarizes the above containment conclusions which Campbell proved.

\begin{figure}[t]\label{fig:1}
$$\includegraphics[scale=.5,angle=0]{rectangles-left} $$
\caption{Rectangles $R_1, R_2, R_3$ contained in Region $A_1$}
\end{figure}


Using Lemma 1 \cite{DMC03}, Campbell showed 

\begin{equation} \label{fund-1}
\left | \frac{f'(z)}{F'(z)}\right | \le 
\left ( \frac{1+G}{1-G} \right )^{\alpha-1} \left [(1-x^2)\frac{C}{B^2} \right ] = {\cal G} \cdot {\cal H} = {\cal L}
\end{equation}

\noindent where

\begin{eqnarray*}
F = F(c,x,a) &=& a + 2c + ac^2 \\
E = E(c,x,a) &=& 1 - c \\
D = D(c,x,a) &=& 1 + ac - x^2(a+c) \\
C = C(c,x,a) &=& |F|(1-x^2) + (x^2-r^2)(1-a^2) \\ 
             &=& |a+2c+ac^2|(1-x^2) + (x^2-r^2)(1-a^2) \\
B = B(c,x,a) &=& |D| - x(1-a)|E| \\
             &=& |1+ac-x^2(a+c)|-x(1-a)|1-c| \\
G = G(c,x,a) &=& x(1-a)\frac{|E|}{|D|}\\
             &=& \frac{x(1-a)|1-c|}{|1+ac-x^2(a+c)|} \\
{\cal G} = {\cal G}(c,x,a) &=&  \left ( \frac{1+G(c,x,a)}{1-G(c,x,a)} \right) ^ {\alpha-1} \\
{\cal H} = {\cal H}(c,x,a) &=& (1-x^2)\frac{C(c,x,a)}{[B(c,x,a)]^2} \\
{\cal L} = {\cal L}(c,x,a) &=& {\cal G}(c,x,a) \cdot {\cal H}(c,x,a) 
\end{eqnarray*}

\begin{figure}[t]\label{fig:2}
$$\includegraphics[scale=.5,angle=0]{rectangles-right} $$
\caption{Rectangles ${\tilde R}_1, {\tilde R}_2, {\tilde R}_3$ contained in Region $A_2$}
\end{figure}

Campbell showed in a separate lemma (Lemma 3 \cite{DMC03}) that ${\cal L}(c,x,a)$ as a function of $\theta$ maximizes at $\theta = 0$, i.e.,  ${\cal L}(r,x,a) \ge {\cal L}(c,x,a)$ for $c = re^{i\theta}$, for $0 \le r \le x \le m(\alpha)$, on the set defined by $Q(a,\alpha) \ge 0$ where 

\begin{equation} \label{poly_Q}
Q(a,\alpha)= -(1+\alpha+2\alpha^2+\alpha^3)a^2 - (2-2\alpha-5\alpha^2-2\alpha^3)a - (\alpha+1). 
\end{equation}

\noindent Finally, Campbell showed that
${\cal L}(r,x,a)$ is an increasing function of $a$ provided that $x \le m(\alpha)$. A calculation shows that ${\cal L}(r,x,1)\equiv 1$.


Let $C_2 = \{(a,\alpha) : Q(c,x,a) = 0\}$ and $A_2 = \{(a,\alpha) : Q(c,x,a) \ge 0\}$. 
For $(a,\alpha) \in A_2$, we have $|f'(z)/F'(z)| \le 1$.  In this notation, Campbell explicitly verified that the rectangles

\begin{eqnarray*}
{\tilde R}_1 = \{(a,\alpha) : 3/20 \le a \le 1, 1.65 \le \alpha \le 2 \}\\
{\tilde R}_2 = \{(a,\alpha) : 1/6 \le a \le 1, 2 \le \alpha \le 3 \}\\ 
{\tilde R}_3 = \{(a,\alpha) : 1/10 \le a \le 1, 3 \le \alpha < \infty \}
\end{eqnarray*} 

\noindent each belong to $A_2$. More explicitly, Campbell showed that the ``left-hand" edges of 
${\tilde R}_1, {\tilde R}_2$ and ${\tilde R}_3$ belong to $A_2$. 

Figure 2 summarizes the above containment conclusions which Campbell proved.  Note, the ``right-hand" edges of the rectanges $R_1, R_2, R_3$ were precisely the ``left-hand" edges of the rectangles ${\tilde R}_1, {\tilde R}_2$ and ${\tilde R}_3$.  Figure 3 summarizes both of the above containment geometries.


Summarizing the above, Campbell proved that subordination $f \prec F$ for $F \in \cU_\alpha$ for $1.65 \le \alpha < \infty$ implies majorization
 $|f'(z)| \le |F'(z)|$ for $|z| \le m(\alpha)$.  

\begin{figure}[t]\label{fig:3}
$$\includegraphics[scale=.5,angle=0]{rectangles} $$
\caption{Rectangles for $\alpha \ge 1.65$}
\end{figure}


%\textbf{See Maplework sheet: implicit plots of curves c1 and c2.mws} 

Campbell conjectured that subordination $f \prec F$ for $F \in \cU_\alpha$ for $1\le \alpha < 1.65$ implies majorization
$|f'(z)| \le |F'(z)|$ for $|z| \le m(\alpha)$.  In \cite{BK}, this conjecture was verified for the case $\alpha = 1$.  For this case, the family $\cU_1$ consists of normalized convex functions.  The proof for the case $\alpha =1$ given in \cite{BK} depended heavily on the convexity of the superordinate function and was not extendable for $\alpha > 1$.  Thus, what remained was to verify Campbell's conjecture for the range $1 < \alpha < 1.65$.


Campbell's arguments, as summarized above, imply, in fact, that for $(a,\alpha) \in A_1 \cup A_2$ that majorization holds even for $\alpha < 1.65$. The boundary curves defining the region $A_1 \cup A_2$ are given implicitly via the bivariate polynomials $p$ and $Q$, as described above, and are difficult to handle analytically.  The preliminary step to verifying Campbell's conjecture consists of identifying two explicit line segments $L_1$ and $L_2$ and analytically verifying that $L_1$ lies in $A_1$ and that $L_2$ lies in $A_2$. See Figure 4.  

\section{Step 0}

\begin{lemma}
Case 1. Let $A_1 = \{(a,\alpha) : p(a,\alpha) \le 1\}$ where $p$ is defined by (\ref{poly_p}).  
Let ${\ds l_1(\alpha) = \frac{\alpha-1}{5}, 1 \le \alpha \le 1.65}$ and 
$L_1 = \{(a,\alpha) : a = l_1(\alpha), 1 \le \alpha \le 1.65\}$.  Then, $L_1 \subset  A_1$. 

Case 2. Let $A_2 = \{(a,\alpha) : Q(a,\alpha) \ge 0\}$ where $Q$ is defined by (\ref{poly_Q}).  
Let ${\ds l_2(\alpha) = \frac{4-2\alpha}{5}, 1 \le \alpha \le 1.65}$ and 
$L_2 = \{(a,\alpha) : a = l_2(\alpha), 1 \le \alpha \le 1.65\}$.  Then, $L_2 \subset A_2$.
\end{lemma}

\textbf{Proof} Case 1. For $1 \le \alpha \le 1.65$, define $k_1(\alpha) = p(l_1(\alpha),\alpha)$.  
Then, we have

$$
{\ds k_1(\alpha) = \frac{((\alpha+1)(\frac{\alpha-1}{5})+1)\left ( \frac{4\alpha+6}{5\alpha}\right)^{\alpha}}{\frac{6\alpha+4}{5}} }
$$

\noindent As a first step, to assure that our arguments will require only finite arithmetic with no numerical approximations, we will make estmates with rational functions whose coefficients are integers.  We will show for $1 \le \alpha \le 1.65$, that 

\begin{equation} \label {ineq_p2}
\left (\frac{4\alpha+6}{5\alpha} \right )^{\alpha} \le p_2(\alpha)
\end{equation}

\noindent where 
${\ds p_2(\alpha) = 2 + \frac{1}{5}(\alpha-1) - \frac{1}{5}(\alpha-1)^2}$.  Then, we will have for $1 \le \alpha \le 1.65$ 

$$
k_1(\alpha) \le k_2(\alpha) = \frac{((\alpha+1)(\frac{\alpha-1}{5})+1)p_2(\alpha) }{\frac{6\alpha+4}{5}} = \frac{n_2(\alpha)}{d_2(\alpha)}.
$$
However, $k_2(\alpha)$ is rational in $\alpha$. A straightforward rearrangement shows that 
$d_2(\alpha) - n_2(\alpha) = (\alpha-1)(\alpha-2)(\alpha^2-6) \ge 0$.  

To verify (\ref{ineq_p2}), we will show 
that 

$$ 
f_2(\alpha) = \log p_2(\alpha) - \alpha \log \frac{4\alpha+6}{5\alpha}  \ge 0
$$

\noindent However, 

$$
f'_2(\alpha) = \frac{7\alpha^2-9\alpha-33}{(\alpha^2-3\alpha-8)(2\alpha+3)}  - \log \frac{4\alpha+6}{5\alpha}
$$

\noindent We will show for $1 \le \alpha \le 1.65$ that

\begin{equation} \label{ineq_p3}
\log \frac{4\alpha+6}{5\alpha} \le p_3(\alpha)
\end{equation}

\noindent where
${\ds p_3(\alpha) = \frac{695}{1000} - \frac{52}{100}(\alpha-1) + \frac{1}{5}(\alpha-1)^2}$.

\noindent Hence, 

$$
f'_2(\alpha) \ge g(\alpha) = \frac{7\alpha^2-9\alpha-33}{(\alpha^2-3\alpha-8)(2\alpha+3)}  - p_3(\alpha)
$$

\noindent However, a Sturm sequence argument, using only finite arithmetic (see \cite{Ja}), shows for $1 \le \alpha \le 1.65$ that $g(\alpha) \ge 0$.  

Finally, to verify (\ref{ineq_p3}), we note that if we let ${\ds p_4(\alpha) = p_3(\alpha) - \log \frac{4\alpha+6}{5\alpha}}$, then
a Sturm sequence argument shows that, $p'_4(\alpha) \ge 0$ and a calculation shows that $p_4(1) > 0$.

Case 2. For $1 \le \alpha \le 1.65$, define $j_1(\alpha) = Q(l_2(\alpha),\alpha)$.  Then, 

$$
j_1(\alpha) = -\frac{4}{25}\alpha^5 -\frac{12}{25}\alpha^4 +\frac{2}{25}\alpha^3 
+\frac{12}{5}\alpha^2 +\frac{7}{5}\alpha -\frac{81}{25}.
$$ 

\noindent A Sturm sequence argument shows that $j_1(\alpha) \ge 0$. $\Box$

\vspace{0.2in}

%\textbf{See Maplework sheet: verify that L1 subset of A1 and L2 subset of A2.mws}

%\textbf{See Maplework sheet: implicit plots of curves c1 and c2.mws}

Campbell's arguments imply for $(a,\alpha)$ to the ``left" of $L_1$ or for $(a,\alpha)$ to the ``right" of $L_2$, that
subordination $f \prec F$ for  $F \in \cU_\alpha$ implies $|f'(z)| \le |F'(z)|$ for $|z| \le m(\alpha)$.
Let $T$ be the critical trapezoid given by

$$T = \{ (a,\alpha) : l_1(\alpha) \le a \le l_2(\alpha), 1 \le \alpha \le 1.65$$

\begin{figure}[t]\label{fig:4}
$$\includegraphics[scale=.5,angle=0]{trapezoid} $$
\caption{Critical Trapezoid}
\end{figure}


\noindent We will show that subordination $f \prec F$ for  $F \in \cU_\alpha$ and for $(a,\alpha) \in T$ implies $|f'(z)| \le |F'(z)|$ for $|z| \le m(\alpha)$.  
More specifically, we will show that subordination $f \prec F$ for  $F \in \cU_\alpha$ and for $(a,\alpha) \in T$ 
implies for $|z| \le m(\alpha)$

\begin{equation} \label{fund-new}
\left | \frac{f'(z)}{F'(z)}\right | \le  \max_{(a,\alpha)\in T}{\cal G}(c,x,a) \cdot \max_{(a,\alpha)\in T}{\cal H}(c,x,a) \le 1
\end{equation}

\noindent That, combined with the observations about Campbell's original arguments, will imply Campbell's conjecture.


\section{Step 1}

Recall that ${\ds {\cal G}(c,x,a) = \left ( \frac{1+G(c,x,a)}{1-G(c,x,a)} \right ) ^{\alpha-1}}$ where  
${\ds G(c,x,a) = \frac{x(1-a)|1-c|}{|1+ac-x^2(a+c)|}}$, $c = re^{i\theta}, 0 \le r \le x \le m(\alpha)$. 
It is easily verified, for $\alpha \ge 1$, that 
${\ds g(y) = \left ( \frac{1+y}{1-y} \right ) ^{\alpha-1}}$  is an non-decreasing function of $y$.  
We will show on $T$ that ${\ds G(c,x,a) \le \frac{6-\alpha}{6+9\alpha}}$.

\vspace{0.2in}

First, we will show that ${ G}(c,x,a) \le { G}(-r,x,a)$, i.e., that ${ G}(c,x,a)$ maximizes at $c = -r$.  
It suffices to show that ${\ds G_1 = \frac{|1-c|^2}{|1+ac-x^2(a+c)|^2}}$ maximizes at $c = -r$.  
Rewriting, we have 
$$
G_1 = \frac{1-2r\cos\theta + r^2}{(1-ax^2)^2 + 2(1-ax^2)(a-x^2)r\cos\theta + (a-x^2)^2r^2} = \frac{\mu_1 + \nu_1 \cos\theta}{\sigma_1 + \tau_1 \cos \theta}.
$$
Hence, 
$$
\frac{\partial G_1}{\partial \theta} = \frac{\mu_1\tau_1 - \nu_1\sigma_1}{(\sigma_1 + \tau_1 \cos\theta)^2}(\sin\theta).
$$
Rewriting, we have
$$
\mu_1\tau_1 - \nu_1\sigma_1 = 2r(1+a)(1-x^2)[1-ax^2 + (a-x^2)r^2]
$$
If $a \ge x^2$, then clearly $\mu_1\tau_1 - \nu_1\sigma_1 \ge 0$.  On the other hand, if $a < x^2$, then
\begin{eqnarray*}
\mu_1\tau_1 - \nu_1\sigma_1 &=& 2r(1+a)(1-x^2)[1-ax^2 + (a-x^2)r^2] \\
                    &\ge& 2r(1+a)(1-x^2)[1-ax^2 + (a-x^2)1] \\
                    &=& 2r (1+a)^2(1-x^2)^2 \\
                    &\ge& 0
\end{eqnarray*}

\vspace{0.2in}


Similarly, by taking derivatives it follows that ${ G}(-r,x,a)$, ${ G}(-x,x,a)$ and ${ G}(-m(\alpha),m(\alpha),a)$ are increasing in $r$, increasing in $x$ and decreasing in $a$, resp.

%Second, we will show that ${\ds {G}(-r,x,a)=\frac{x(1-a)(1+r)}{1-ar-x^2(a-r)}}$ is increasing in $r$.  
%It suffices to show that ${\ds G_2 = \frac{(1+r)}{1-ar-x^2(a-r)} = \frac{\mu_2 + \nu_2 r}{\sigma_2 + \tau_2 r}}$ is increasing %in $r$.  Hence,
%$
%{\ds \frac{\partial G_2}{\partial \theta} = \frac{-\mu_2\tau_2 + \nu_2\sigma_2}{(\sigma_2 + \tau_2 r)^2}}.
%$
%Here, $-\mu_2\tau_2 + \nu_2\sigma_2 = 1(1-ax^2) + 1(a-x^2) = (1-x^2)(1+a)\ge 0$.

%\vspace{0.2in}

%Third, we will show that ${\ds { G}(-x,x,a) = \frac{x(1-a)(1+x)}{1-ax-x^2(a-x)}}$ is increasing in $r$.  
%It suffices to show that 
%$$
%G_3 = \frac{x(1+x)}{1-ax-x^2(a-x)} = \frac{x(1+x)}{(1+x)(1-x-ax+x^2)} = \frac{x}{1-x-ax+x^2}
%$$
%is increasing in $x$. But, ${\ds \frac{\partial G_3}{\partial \theta} = \frac{1-x^2}{(1-x-ax+x^2)^2} } \ge 0$.

%\vspace{0.2in}

%Fourth, we will show that ${ G}(-m(\alpha),m(\alpha),a)$ is decreasing in $a$.  Let $m = m(\alpha)$.  Then,
%${\ds G(-m,m,a) = \frac{m(1-a)(1+m)}{1-am-m^2(a-m)}} $.  It suffices to show that
%\begin{eqnarray*}
%G_4 &=& \frac{(1+m)(1-a)}{1-am-m^2(a-m)} = \frac{(1+m)(1-a)}{(1+m)(1-m-am+m^2)} \\
%    &=& \frac{1-a}{1-m-am+m^2}
%\end{eqnarray*}
%is decreasing in $a$. Hence, 
%$
%{\ds \frac{\partial G_4}{\partial a} = -\frac{(1-m)^2}{(1-m-am+m^2)^2} \le 0 }.
%$ 

\vspace{0.2in}

Thus, on $T$ we have ${ G}(c,x,a) \le { G}(-m(\alpha),m(\alpha),l_1(\alpha))$.  It is a straightforward calculation, using the relationship that $1 + m(\alpha)^2 = 2(\alpha+1)m(\alpha)$, that 
${\ds {G}(-m(\alpha),m(\alpha),l_1(\alpha))= \frac{6-\alpha}{6+9\alpha}}$.  


\section{Step 2}

Recall that ${\ds {\cal G}(c,x,a) = \left ( \frac{1+G(c,x,a)}{1-G(c,x,a)} \right ) ^{\alpha-1}}$ where  
${\ds G(c,x,a) = \frac{x(1-a)|1-c|}{|1+ac-x^2(a+c)|}}$, $c = re^{i\theta}, 0 \le r \le x \le m(\alpha)$. 
In Step 1, we noted, for $\alpha \ge 1$, that 
${\ds g(y) = \left ( \frac{1+y}{1-y} \right ) ^{\alpha-1}}$  is a non-decreasing function of $y$.  
Further, in Step 1, we showed on $T$ that ${\ds G(c,x,a) \le \frac{6-\alpha}{6+9\alpha}}$.

Again, to assure that our arguments will require only finite arithmetic with no numerical approximations, we will make estmates with rational functions whose coefficients are integers. Let ${\ds l(y) = 1 +\frac{21}{10}(\alpha-1)(1+\frac{\alpha-1}{4})y}$. We will show that ${\ds g(\frac{6-\alpha}{6+9\alpha}) \le l(\frac{6-\alpha}{6+9\alpha})}$.  Hence, on $T$ we will have that ${\ds {\cal G}(c,x,a) \le l(\frac{6-\alpha}{6+9\alpha}})$.

Let ${\ds h(\alpha) = l(\frac{6-\alpha}{6+9\alpha}) - g(\frac{6-\alpha}{6+9\alpha}) }$.  We can rewrite $h$ as 
${\ds h(\alpha) = \frac{n(\alpha)}{40(2+3\alpha)} }$ where
$$
n(\alpha) = -46 + 225\alpha + 28\alpha^2 - 7\alpha^3 -40(2+3\alpha)\left( \frac{2(3+2\alpha)}{5\alpha} \right)^{\alpha-1}.
$$
We will show explicitly for $1 \le \alpha \le 1.65$ that $n$ is an non-negative increasing function of $\alpha$.
Let ${\ds \beta = \left( \frac{2(3+2\alpha)}{5\alpha} \right)^{\alpha-1} }$ and ${\ds \gamma = \log \frac{2(3+2\alpha)}{5\alpha} }$.  Then, we can write
$$
n'(\alpha) = \frac{n_1(\alpha)}{\alpha(3+2\alpha)} =\frac{A_0  + \beta(B_0  +\gamma C_0)}{\alpha(3+2\alpha)}
$$
where
\begin{eqnarray*}
A_0 &=& 675\alpha + 618\alpha^2 +49\alpha^3 -42\alpha^4 \ge 0\\
B_0 &=& -240 -480\alpha+120\alpha^2 \le 0\\
C_0 &=& -240\alpha-520\alpha^2-240\alpha^3 \le 0.
\end{eqnarray*}
It is sufficient to show that $n_1$ is non-negative.

In separate lemmas we will show for $1 \le \alpha \le 1.65$ the following upper bounds for $\beta$ and $\gamma$ hold
\begin{eqnarray} \label{bnd_beta}
\beta &\le& \beta_1 = 1+\frac{694}{1000}(\alpha-1)-\frac{361}{1000}(\alpha-1)^2+\frac{69}{1000}(\alpha-1)^3 \\ \label{bnd_gamma}
\gamma &\le& \gamma_1 = \frac{6932}{10000}-\frac{5992}{10000}(\alpha-1)+\frac{42}{100}(\alpha-1)^2-\frac{2795}{10000}(\alpha-1)^3 \\ \nonumber
&+&\frac{111}{1000}(\alpha-1)^4
\end{eqnarray}
Using the upper bounds for $\beta$ and $\gamma$ in (\ref{bnd_beta}) and (\ref{bnd_gamma}), we have
$$
n_1(\alpha) \ge n_2(\alpha) = A_0 + \beta_1(B_0 + \gamma_1 C_0)
$$
However, a Sturm sequence argument shows for $1 \le \alpha \le 1.65$ that $n_2(\alpha) \ge 0$.

\vspace{0.2in}

\section{Step 3}

a) Let $(a,\alpha) \in T$ be fixed and let $x \le m$ be fixed.  For $c = re^{i\theta}$ we show that ${\cal H}(c,x,a)$ does not have a local extreme value on the open 
set $O_r = \{(r,\theta) : 0 < r < x, 0 < \theta < \pi \}$.  Specifically, if ${\cal H}(c,x,a)$ had a local extreme value, say at $(r,\theta)$, then
we would have 

\begin{eqnarray} \label{partial_r_theta}
\frac{\partial {\cal H}}{\partial r}{(r,\theta)} &=& 0 \\ \nonumber
\frac{\partial {\cal H}}{\partial \theta}{(r,\theta)} &=& 0
\end{eqnarray}

However, (\ref{partial_r_theta}) implies that 

\begin{eqnarray*} 
B(r,\theta)\frac{\partial {C}}{\partial r}{(r,\theta)} - 2C(r,\theta)\frac{\partial {B}}{\partial r}{(r,\theta)}&=& 0 \\ \nonumber
B(r,\theta)\frac{\partial {C}}{\partial \theta}{(r,\theta)} - 2C(r,\theta)\frac{\partial {B}}{\partial \theta}{(r,\theta)}&=& 0 
\end{eqnarray*}

which in turn implies that 

\begin{equation} \label{det-r}
\Delta_r = \left |
\begin{array}{ll}
\frac{\partial {C}}{\partial r}{(r,\theta)} & \frac{\partial {B}}{\partial r}{(r,\theta)} \\
& \\
\frac{\partial {C}}{\partial \theta}{(r,\theta)} & \frac{\partial {B}}{\partial \theta}{(r,\theta)} \\
\end{array}
\right | = 0
\end{equation}

We will show that the determinant in (\ref{det-r}) is, however, non-zero.  Specifically, computing we have
$$
\Delta_r = 2 r^2 \sin \theta (1-a^2) \frac{N_r}{D_r} 
$$
where
\begin{eqnarray*}
N_r &=& (p_r + q_r |F|) + (P_r + Q_r |F|)G \\
D_r &=& |E| |F|
\end{eqnarray*}
and $p_r, q_r, P_r$ and $Q_r$ are polynomials, with rational coefficients, in $r, \cos \theta, x$ and $a$.

It suffices to show that $N_r$ is negative. Since $N_r$ is linear in $G$ and, from Step 1, we have that 
${\ds 0 \le G \le \frac{6-\alpha}{6+9\alpha} }$, then $ng0 \le N_r \le ng1$
where
\begin{eqnarray*}
ng0 &=& (p_r + q_r  |F|) + (P_r + Q_r |F|)0 \\
ng1 &=& (p_r + q_r  |F|) + (P_r + Q_r |F|)\frac{6-\alpha}{6+9\alpha} \\
\end{eqnarray*}.

Further, since $ng0$ and $ng1$ are each linear in $|F|$ and $0 \le |F| \le a + 2x + ax^2$, then we have
$ng00 \le ng0 \le ng01$ and $ng10 \le ng1 \le ng11$
where
\begin{eqnarray*}
ng00 &=& (p_r + q_r 0) + (P_r + Q_r 0)0  = p_r\\
ng01 &=& (p_r + q_r (a + 2x + ax^2)) + (P_r + Q_r (a + 2x + ax^2))0  = p_r + q_r (a + 2x + ax^2)) \\
ng10 &=& (p_r + q_r 0) + (P_r + Q_r 0)\frac{6-\alpha}{6+9\alpha} = p_r + P_r\frac{6-\alpha}{6+9\alpha} \\
ng11 &=& (p_r + q_r (a + 2x + ax^2)) + (P_r + Q_r (a + 2x + ax^2))\frac{6-\alpha}{6+9\alpha} \\
\end{eqnarray*}.

In a lengthy, finite arithmetic argument, using exact calcluations, in the Maple worksheet {\em campbell-maple-1.mws} we explicitly verify that each of the factors $ng00, ng01, ng10, ng11$ are negative.

b) Since ${\cal H}(c,x,a)$ does not have a local extreme value on the open 
set $O_r = \{(r,\theta) : 0 < r < x, 0 < \theta < \pi \}$, then any extreme value of ${\cal H}(c,x,a)$ must occur on either 

b1) the set $E(r)^{+}= \{(r,0) : 0 \le r \le x \}$

b2) the set $E(r)^{-} =\{(-r,0) : 0 \le r \le x \}$

b3) the set $E(x,\theta) = \{(x,\theta) : 0 \le \theta \le \pi \}$

\vspace{0.1in}

For $r \in E(r)^{+}$, $a + 2r + ar^2$ does not change sign.  Hence, we explicitly have
$$
{\cal H}(r,x,a) = (1-x^2)\frac{(a+2r+ar^2)(1-x^2)+(x^2-r^2)(1-a^2)}{( 1 + ar-x^2(a+r) - x(1-a)(1-r) )^2} .
$$
A straightforward exercise shows that ${\cal H}(r,x,a)$ is increasing in both $r$ and $a$.  Hence, 
on $E(r)^{+}$ the maximum value of ${\cal H}(r,x,a)$ is ${\cal H}(x,x,l_2(\alpha))$ 
%(See the Maple worksheet: {\em verify step 3 b on set E(r)1.mws})

\vspace{0.1in}

On the set $E(r)^{-}$, $a -2r + ar^2$ changes sign whenever $a = \frac{2r}{1+r^2}$.  Hence, we have,
generically, 
$$
{\cal H}(r,x,a) = (1-x^2)\frac{|a-2r+ar^2|(1-x^2)+(x^2-r^2)(1-a^2)}{( 1 - ar-x^2(a-r) - x(1-a)(1+r) )^2} .
$$
Divide the domain space $D = \{(a,r) : 0 \le r \le x, l_1(\alpha) \le a \le l_2(\alpha)\}$ into subsets on which $a -2r + ar^2$ has constant sign, i.e., $D = D_1 \cup D_2$ where $D_1 = \{ (a,r) \in D : a -2r + ar^2 \ge 0 \}$ and 
$D_2 = \{ (a,r) \in D : a -2r + ar^2 \le 0 \}$.  On $D_1$ it is a straightforward exercise to show that ${\cal H}(-r,x,a)$ is increasing in $a$ and decreasing in $r$.  Hence, ${\cal H}(-r,x,a) \le {\cal }H(0,x,l_2(\alpha)$.  However, by part b1) 
${\cal H}(0,x,l_2(\alpha)) \le {\cal H}(x,x,l_2(\alpha))$.  On $D_2$ it follows that ${\cal H}(-r,x,a)$ is decreasing in $a$ and increasing in $r$.  Hence, ${\cal H}(-r,x,a) \le {\cal }H(-x,x,l_1(\alpha)$.  Consequently, on $D$ we have
$$
{\cal H}(-r,x,a) \le \max \{{\cal H}(x,x,l_2(\alpha)),{\cal H}(-x,x,l_1(\alpha))\}
$$
%(See the Maple worksheet: {\em verify step 3 b on set E(r)2.mws})

Note, each of the potential extreme values for the cases b1) and b2) occur on the set $E(x,\theta)$.

\vspace{0.1in}

c) Consider the set $E(x,\theta)$.  Let $(a,\alpha) \in T$ be fixed.  For $c = xe^{i\theta}$ 
we will show that ${\cal H}(c,x,a)$ 
does not have a local extreme value on the open set $O_x = \{(x,\theta) : 0 < x < m, 0 < \theta < \pi \}$.  
Specifically, if ${\cal H}(c,x,a)$ had a local extreme value, say at $(x,\theta)$, then
we would have 

\begin{eqnarray} \label{partial_x_theta}
\frac{\partial {\cal H}}{\partial x}{(x,\theta)} &=& 0 \\ \nonumber
\frac{\partial {\cal H}}{\partial \theta}{(x,\theta)} &=& 0
\end{eqnarray}

However, (\ref{partial_x_theta}) implies that 

\begin{eqnarray*} 
B(x,\theta)\frac{\partial {C}}{\partial x}{(x,\theta)} - 2C(x,\theta)\frac{\partial {B}}{\partial x}{(x,\theta)}&=& 0 \\ \nonumber
B(x,\theta)\frac{\partial {C}}{\partial \theta}{(x,\theta)} - 2C(x,\theta)\frac{\partial {B}}{\partial \theta}{(x,\theta)}&=& 0 
\end{eqnarray*}

which in turn implies that 

\begin{equation} \label{det-x}
\Delta_x = 
\left |
\begin{array}{ll}
\frac{\partial {C}}{\partial x}{(x,\theta)} & \frac{\partial {B}}{\partial x}{(x,\theta)} \\
& \\
\frac{\partial {C}}{\partial \theta}{(x,\theta)} & \frac{\partial {B}}{\partial \theta}{(x,\theta)} \\
\end{array}
\right | = 0
\end{equation}

\vspace{0.1in}

We will show that the determinant in (\ref{det-x}) is, however, non-zero.  Specifically, computing we have
$$
\Delta_x = 2 x \sin \theta (1-x^2) \frac{N_x}{D_x} 
$$
where $N_x = P_x + Q_x G, D_x = |E| |F|$
and $P_x, Q_x$ are polynomials, with rational coefficients, in $r, \cos \theta, x$ and $a$.

It suffices to show that $N_x$ is negative. Since $N_x$ is linear in $G$ and, from Step 1, we have that 
${\ds 0 \le G \le \frac{6-\alpha}{6+9\alpha} }$, then $nG0 \le N_r \le nG1$
where
\begin{eqnarray*}
nG0 &=& P_x + Q_x 0  = P_x\\
nG1 &=& P_x + Q_x \frac{6-\alpha}{6+9\alpha} \\
\end{eqnarray*}.

In a lengthy, finite arithmetic argument, using exact calcluations, in the Maple worksheet {\em campbell-maple-2.mws} we explicitly verify that each of the factors $nG0, nG1$ are negative.

d) Since ${\cal H}(c,x,a)$ does not have a local extreme value on the open 
set $O_x = \{(x,\theta) : 0 < x < m, 0 < \theta < \pi \}$, then any extreme value of ${\cal H}(c,x,a)$ must occur on either 

d1) the set $E(x)^{+} = \{(x,0) : 0 \le x \le m \}$

d2) the set $E(x)^{-} =\{(-x,0) : 0 \le x \le m \}$

d3) the set $E(m,\theta) = \{(m,\theta) : 0 \le \theta \le \pi \}$

\vspace{0.1in}

\vspace{0.1in}

On the set $E(x)^{+}$, $a + 2x + ax^2$ does not change sign.  Hence, we explicitly have
$$
{\cal H}(x,x,a) = \frac{(a+2x+ax^2)(1-x^2)^2}{( 1 + ax-x^2(a+x) - x(1-a)(1-x) )^2} .
$$
A straightforward exercise shows that ${\cal H}(x,x,a)$ is increasing in both $x$ and $a$.  Hence, 
on $E(x)^{+}$ the maximum value of ${\cal H}(x,x,a)$ is ${\cal H}(m,m,l_2(\alpha))$ 
%(See the Maple worksheet: {\em verify step 3 d on set E(x)1.mws})

\vspace{0.1in}

On the set $E(x)^{-}$, $a -2x + ax^2$ changes sign whenever $a = \frac{2x}{1+x^2}$.  Hence, we have,
generically, 
$$
{\cal H}(x,x,a) = \frac{|a-2x+ax^2|(1-x^2)^2}{( 1 - ax-x^2(a-x) - x(1-a)(1+x) )^2} .
$$
Divide the domain space $D = \{(a,x) : 0 \le x \le m, l_1(\alpha) \le a \le l_2(\alpha)\}$ into subsets on which $a -2x + ax^2$ has constant sign, i.e., $D = D_1 \cup D_2$ where $D_1 = \{ (a,r) \in D : a -2x + ax^2 \ge 0 \}$ and 
$D_2 = \{ (a,x) \in D : a -2x + ax^2 \le 0 \}$.  On $D_1$ it is a straightforward exercise to show that ${\cal H}(-x,x,a)$ is increasing in $a$ and decreasing in $x$.  Hence, ${\cal H}(-x,x,a) \le {\cal }H(0,0,l_2(\alpha)$.  However, by part d1) 
${\cal H}(0,0,l_2(\alpha)) \le {\cal H}(m,m,l_2(\alpha))$.  On $D_2$ it follows that ${\cal H}(-x,x,a)$ is decreasing in $a$ and increasing in $x$.  Hence, ${\cal H}(-x,x,a) \le {\cal }H(-m,m,l_1(\alpha)$.  Consequently, on $D$ we have
$$
{\cal H}(-x,x,a) \le \max \{{\cal H}(m,m,l_2(\alpha)),{\cal H}(-m,m,l_1(\alpha))\}
$$
%(See the Maple worksheet: {\em verify step 3 d on set E(x)2.mws})

Note, each of the potential extreme values for the cases b1) and b2) occur on the set $E(m,\theta)$.

Further note, using the relationship $1 + m^2 = 2(\alpha+1)m$, it can be explicity shown
\begin{eqnarray} \label{explicit}
{\cal H}(m,m,l_2(\alpha)) &=& \frac{(\alpha+1)l_2(\alpha)+1}{(\alpha+l_2(\alpha)+1)^2}(\alpha+2) \\
{\cal H}(-m,m,l_1(\alpha)) &=& \frac{1-l_1(\alpha)(\alpha+1)}{\alpha} \nonumber
\end{eqnarray}


Finally, consider the set $E(m,\theta) = \{(m,\theta) : 0 \le \theta \le \pi \}$. 
Let $(a,\alpha) \in T$ and $x = m(\alpha)$.  For $c = xe^{i\theta}$, we show that 
\begin{equation} \label{last}
{\cal H}(c,x,a) \le h_3(\alpha) = 1 - \frac{4}{5}(\alpha-1) + \frac{13}{10}(\alpha-1)^2 - \frac{13}{10}(\alpha-1)^3
\end{equation}  
We note that $h_3(\alpha)$ is an upper bound for the sharp bounds on ${\cal H}(c,x,a)$ on the sets $E_1$ and $E_2$ given in (\ref{explicit}). In a lengthy, finite arithmetic argument, in the Maple worksheet {\em campbell-maple-3.mws}, the inequality (\ref{last}) is verified.

\vspace{0.1in}

Secondary lemmas required for this last step are:


\begin{lemma}[A-3]  \label{lem:poly3}
Let $p = p(x_1,x_2,x_3)$ be a polynomial with rational coefficients defined on the region 
$R = \{ (x_1,x_2,x_3) : a_1 \le x_1 \le b_1,\ \ a_2 \le x_2 \le b_2,\ \ a_3 \le x_3 \le b_3\}$, where $a_j,\ b_j$ are 
all rational numbers. Let $M$ be an upper bound for  ${\ds |\frac{\partial p}{\partial x_j}(x_1,x_2,x_3)|}$ on $R$, $1 \le j \le 3$.   
Let $\delta > 0 $ and suppose $N_1, N_2, N_3$ are chosen so that
$\Delta_j = (b_j-a_j)/N_j \le \delta$, $1 \le j \le 3$.  Let $L$ be the lattice
$L = L(N_1,N_2,N_3) = \{ (a_1 + i\Delta_1,a_2 + j\Delta_2,a_3 + k\Delta_3)\  :\  0 \le i \le N_1,\ \   0 \le j \le N_2,
\ \  0 \le k \le N_3\}$.
Let 
$$m = \min_{(x_1,x_2,x_2) \in L} p(x_1,x_2,x_3) . $$  
If $m \ge \frac{3}{2}M\delta$, then $p(x_1,x_2,x_3) \ge 0$ on $R$.
\end{lemma}
\textbf{Proof}
Let $(x_1,x_2,x_3) \in R$.  Then, there exists $(x_1^0,x_2^0,x_3^0) \in L$ such that ${\ds {\rm dist}( (x_1^0,x_2^0,x_3^00),(x_1,x_2,x_3) ) \le \frac{\sqrt{3}}{2}\delta}$.  Without loss of generality we can suppose that 
$x_1 \ge x_1^0, x_2 \ge x_2^0, x_3 \ge x_3^0$.
Define 
$$p(t) = p_{\theta,\phi}(t) = p(x_1^0+t\cos(\theta)\sin(\phi),x_2^0+t\sin(\theta)\sin(\phi),x_3^0+t\cos(\phi))$$ 
$0 < t \le \frac{\sqrt{3}}{2}\delta, 0 < \theta < \pi/2, 0 < \phi < \pi/2$.  Then, $\left | p'(t) \right | \le M\sqrt{3}$ for $0 \le t \le \frac{\sqrt{3}}{2}\delta$.
Hence, $p(x_1,x_2,x_3) = p_{\theta,\phi}(t) \ge p_{\theta,\phi}(0) - M\sqrt{3}\, t = p(x_1^0,x_2^0,x_3^0) - M \sqrt{3}\ t \ge m - \frac{3}{2}M \delta \ge 0$. $\Box$

\vspace{0.1in}

\begin{lemma}[B-1] \label{lem:est1}
Let $p = p(x)$ be a polynomial with rational coefficients defined on the interval $R_1(d) = \{ x : 0 < x < d\}$ such
that $p(0) = 0$.  Let $M_1 = p'(0) > 0$ and let $-N_1$ be a lower bound for $p''(x)$ on $R_1(d)$ where $N_1 \ge 0$.  Then, $p(x) > 0$ on $R_1(\delta_1)$ where $\delta_1 = \min\{\frac{2 M_1}{N_1},d\}$.
\end{lemma}
\textbf{Proof}
The hypotheses imply, using Taylor's remainder theorem, that 

\noindent ${\ds p(x) = M_1 x + p''(c)\frac{x^2}{2} }$ for some $c$ such that $0 < c < x$.  Hence, we have ${\ds \frac{p(x)}{x} > M_1 - \frac{N_1 x}{2} }$.  $\Box$.

%Since $p'(0) = M_1$ and since $p''(x) \ge -N_1$ on $R_1(d)$, then $p'(x) > 0$ on $R_1(\delta)$.  Hence, $p$ is strictly %increasing on $R_1(\delta)$ which implies since $p(0) = 0$ that $p(x) > 0$ on $R_1(\delta)$.


\begin{lemma}[B-3] \label{lem:est3}
Let $p = p(x_1,x_2,x_3)$ be a polynomial with rational coefficients defined on the region $R_3(d) = \{ (x_1,x_2,x_3) : 0 < x_1 < d\ , 0 < x_2 < d\ ,0 < x_3 < d\}$ such that $p(0,0,0) = 0$.  
Let 
$$M_3 = \min\{p_1(0,0,0), p_2(0,0,0), p_3(0,0,0)\} > 0$$
where $p_j$ denotes the partial of $p$ with respect to $x_j, j = 1, 2, 3$.  
Let $-N_3$ be a lower bound for $p_{ij}(x_1,x_2,x_3)$ on $R_3(d)$, $ 1 \le i,j \le 3$,  where $N_3 \ge 0$.  Then,
$p(x_1,x_2,x_3) > 0$ on $R_3(\delta_3)$ where $\delta_3 = \min\{2 M_3/(3N_3),d\}$.
\end{lemma}
\textbf{Proof}
Let $B_3(d) = \{(x_1,x_2,x_3) \in R_3(d) : ||(x,y,z)|| < d\}$.  Note that $R_3(d/\sqrt{3}) \subset B_3(d)$.  On $B_3(d)$ we can write 
$$p(x_1,x_2,x_3) = p(t\cos(\theta)\sin(\phi),t\sin(\theta)\sin(\phi),t\cos(\phi)) = p_{\theta,\phi}(t) = p(t)$$ 
where $0 < t < d, 0 < \theta < \pi/2, 0 < \phi < \pi/2$.  Apply Lemma B-1 to $p$ on $R_1(d)$ with $M_1 = M_3$ and $N_1 = \sqrt{3}N_3$. Consequently, $p(t) > 0$ on $R_1(\delta_1)$ where $\delta_1 = 2M_3/(\sqrt{3} N_3)$.  Hence, $p(x_1,x_2,x_3) > 0$ on $R_3(\delta_3/\sqrt{3})$. $\Box$

\vspace{0.2in}

\section{Step 4}

Recall in Step 1 and Step 2, we showed that 
$$
{\cal G}(c,x,a) \le l(\frac{6-\alpha}{6+9\alpha}) = g_3(\alpha) =
1 + \frac{21}{10}(\alpha-1)(1 + \frac{\alpha-1}{4})\frac{6-\alpha}{6+9\alpha}.
$$
Further, in Step 3 we showed 
$$
{\cal H}(c,x,a) \le h_3(\alpha) =  1 - \frac{4}{5}(\alpha-1) + \frac{13}{10}(\alpha-1)^2 - \frac{13}{10}(\alpha-1)^3.
$$  

It remains to show for $1 \le \alpha \le 1.65$ that $g_3(\alpha)\cdot h_3(\alpha) \le 1$.  It is sufficient to show
that $k_3(\alpha) \ge 0$ where 
\begin{eqnarray*}
k_3(\alpha) &=& 400(6+9\alpha)(1-g_3(\alpha)\cdot h_3(\alpha)) \\
            &=& 8472-36174\alpha+52755\alpha^2-29838\alpha^3+2874\alpha^4+2184\alpha^5-273\alpha^6
\end{eqnarray*}
However, a Sturm sequence argument verifies that $k_3(\alpha) \ge 0$.

\section{Appendix}

Verification of upper bounds given in Step 2 for $1 \le \alpha \le 1.65$ for $\beta$ and $\gamma$.

Case $\gamma$.
\begin{eqnarray*} 
\gamma &=& \log \frac{2(3+2\alpha)}{5\alpha} \\ &\le& \frac{6932}{10000}-\frac{5992}{10000}(\alpha-1)+\frac{42}{100}(\alpha-1)^2-\frac{2795}{10000}(\alpha-1)^3 + \frac{111}{1000}(\alpha-1)^4 = \gamma_1
\end{eqnarray*}

We will show that $\gamma_2 = \gamma_1 - \gamma$ is a non-negative increasing function.  We have
\begin{eqnarray*}
\frac{d \gamma_2(\alpha)}{d\alpha} &=& \frac{N_1(\alpha)}{10000\alpha(3+2\alpha)} \\ 
      &=& \frac{30000-81651\alpha+61036\alpha^2+11865\alpha^3-30090\alpha^4+8880\alpha^5}{10000\alpha(3+2\alpha)}
\end{eqnarray*}
A Sturm sequence argument shows that $N_1$ is non-negative.

\vspace{0.2in}

Case $\beta$.  
\begin{eqnarray*} 
\beta = \left( \frac{2(3+2\alpha)}{5\alpha} \right)^{\alpha-1} &\le&  1+\frac{694}{1000}(\alpha-1)-\frac{361}{1000}(\alpha-1)^2+\frac{69}{1000}(\alpha-1)^3 = \beta_1
\end{eqnarray*}

It suffices to show that $\beta_2 = \log \beta_1 - \log \beta \ge 0$.  We will show for $1 \le \alpha \le 1.65$ 
that $\beta_2$ is a non-negative increasing function of $\alpha$.  
We have 
\begin{eqnarray} \label{denom_beta}
\frac{d \beta_2(\alpha)}{d\alpha} &=& \frac{N_2(\alpha)}{\alpha((-124+1623\alpha-568\alpha^2+69\alpha^3))(3+2\alpha)}
\end{eqnarray}
where
\begin{eqnarray*} 
N_2(\alpha) &=& 372-372\alpha+6411\alpha^2-3562\alpha^3+621\alpha^4 \\
            &+& (372\alpha-4621\alpha^2-1542\alpha^3+929\alpha^4-138\alpha^5)\gamma
\end{eqnarray*}
A Sturm sequence argument shows for $1 \le \alpha \le 1.65$ that the denominator of (\ref{denom_beta}) is positive.  
We have then, since $\gamma \le \gamma_1$ that $N_2(\alpha) \ge N_3(\alpha)$ where
\begin{eqnarray*} 
N_3(\alpha) &=& 372-372\alpha+6411\alpha^2-3562\alpha^3+621\alpha^4 \\
            &+& (372\alpha-4621\alpha^2-1542\alpha^3+929\alpha^4-138\alpha^5)\gamma_1
\end{eqnarray*}
A Sturm sequence argument shows for $1 \le \alpha \le 1.65$ that $N_3(\alpha) \ge 0$.


\vspace{0.1in}

%\bibliographystyle{amsplain} \bibliography{hc,roger,eqcp,dis}

\def\cprime{$'$}
\providecommand{\bysame}{\leavevmode\hbox to3em{\hrulefill}\thinspace}
\providecommand{\MR}{\relax\ifhmode\unskip\space\fi MR }
% \MRhref is called by the amsart/book/proc definition of \MR.
\providecommand{\MRhref}[2]{%
  \href{http://www.ams.org/mathscinet-getitem?mr=#1}{#2}
}
\providecommand{\href}[2]{#2}


\begin{thebibliography}{10000}

\bibitem{BK}
Roger W. Barnard and Charles Kellogg, \emph{On Campbell's conjecture on the radius of majorization of functions subordinate to convex functions}, Rocky Mountain J. of Mathematics \textbf{14} (1984), 331-339. 

\bibitem{DMC01}
Douglas M. Campbell, \emph{Majorization-subordination theorems for locally univalent functions. I}, Bulletin
  of the American Mathematical Society \textbf{78} (1972), 535-538. \MR{45:8817}.
  
\bibitem{DMC02}
Douglas M. Campbell, \emph{Majorization-subordination theorems for locally univalent functions. II}, Canadian 
  Journal of the American Mathematical Society \textbf{25} (1973), 420-425. 
  
\bibitem{DMC03}
Douglas M. Campbell, \emph{Majorization-subordination theorems for locally univalent functions. III}, Transactions
  of the American Mathematical Society \textbf{198} (1974), 297-306. 
  
\bibitem{Ja} N. Jacobson, 
\textit{Basic Algebra. I.} Second Edition.  W. H. Freeman and Company, New York, 1985. % MR0780184 (86d:00001)

\bibitem{Pom01}
Christian Pommerenke, \emph{Linear-invariante Familien analytishcer Functionen. I}, Mathematische Annalen \textbf{155} 
  (1964), 108-154. \MR{29:6007}.
  
\end{thebibliography}


\end{document}
